%% ACM Transactions on Software Engineering and Methodology (TOSEM) — manuscript submission
%% Compile: latexmk -pdf main
%% Or:      pdflatex main && bibtex main && pdflatex main && pdflatex main

\documentclass[manuscript,screen,review]{acmart}

%% Journal target
\acmJournal{TOSEM}

%% ── ACM Metadata ────────────────────────────────────────────────────────────
\title{Context-Scoped AI Agents in Integrated Developer Workspaces: A Between-Subjects Study of Information Breadth and Task Performance}

%% Authors — replace with real names, affiliations, ORCIDs
\author{Author One}
\authornote{Corresponding author.}
\email{author1@institution.edu}
\orcid{0000-0000-0000-0000}
\affiliation{%
  \institution{Institution Name}
  \streetaddress{Street Address}
  \city{City}
  \state{State}
  \country{Country}
  \postcode{00000}
}

\author{Author Two}
\email{author2@institution.edu}
\affiliation{%
  \institution{Institution Name}
  \city{City}
  \country{Country}
}

%% ── Abstract ────────────────────────────────────────────────────────────────
\begin{abstract}
Software development is fragmented across issue trackers, team chat, and code
hosts, and AI assistants embedded in that workflow face an analogous
fragmentation: they can only act on what they are allowed to see. We
operationalise this as a \emph{context ladder}---a three-rung experimental
manipulation of how much of a shared workspace an AI agent can observe and act
on: chat only (\CI), chat and task board (\CII), or chat, board, and GitHub
(\CIII)---and measure the downstream effect on user task performance, cognitive
load, and trust across a population spanning both novice and experienced
programmers.

We present \Weld, an open-source integrated developer workspace
(SvelteKit\,+\,Rust/Axum) combining a Kanban board, team chat, and mocked
GitHub repository data in one product, and report a hybrid user study (target
$N$~=~20--30) recruiting from two complementary populations: university CS
students (undergraduate and graduate, representing less-experienced developers)
and professional software developers ($\geq$1\,year of industry experience).
Every participant first completes two manual no-agent baseline tasks---
establishing a per-person speed and workload anchor that absorbs individual
skill differences---then is randomly assigned to one agent condition for six
agent-assisted tasks spanning comprehension, single-surface, and cross-surface
coordination types. Including both populations allows us to examine whether the
benefit of context breadth varies with developer experience, and whether novice
and experienced programmers differ in their reliance on and trust of an AI agent
with varying information access. Primary analysis uses baseline-normalised task
time via a mixed-effects model testing the condition~$\times$ task-type
interaction; secondary measures include \NASATLX, System Usability Scale, a
five-item integration-perception scale, and an eight-item AI trust scale.

Context scoping is enforced at the server in code---not in prompts---by the
\ContextGate module, which records a per-turn \texttt{agent\_context\_sent}
audit snapshot proving, from the log alone, exactly what information each agent
invocation received. This architectural choice makes the experimental
manipulation verifiable rather than merely assumed. We contribute the
context-ladder study design, the mixed-experience participant protocol, and the
\ContextGate mechanism as replicable instruments for empirical research on AI
agents in multi-surface collaborative tools.
\end{abstract}

%% ── CCS Concepts ────────────────────────────────────────────────────────────
%% Generate from https://dl.acm.org/ccs
\begin{CCSXML}
<ccs2012>
  <concept>
    <concept_id>10011007.10011006.10011073</concept_id>
    <concept_desc>Software and its engineering~Software development techniques</concept_desc>
    <concept_significance>500</concept_significance>
  </concept>
  <concept>
    <concept_id>10003120.10003121.10003122</concept_id>
    <concept_desc>Human-centered computing~HCI theory, concepts and models</concept_desc>
    <concept_significance>500</concept_significance>
  </concept>
  <concept>
    <concept_id>10003120.10003121.10003124</concept_id>
    <concept_desc>Human-centered computing~Empirical studies in HCI</concept_desc>
    <concept_significance>500</concept_significance>
  </concept>
  <concept>
    <concept_id>10010147.10010178</concept_id>
    <concept_desc>Computing methodologies~Artificial intelligence</concept_desc>
    <concept_significance>300</concept_significance>
  </concept>
</ccs2012>
\end{CCSXML}

\ccsdesc[500]{Software and its engineering~Software development techniques}
\ccsdesc[500]{Human-centered computing~HCI theory, concepts and models}
\ccsdesc[500]{Human-centered computing~Empirical studies in HCI}
\ccsdesc[300]{Computing methodologies~Artificial intelligence}

%% ── Keywords ────────────────────────────────────────────────────────────────
\keywords{AI agents, developer tools, integrated workspaces, context scoping,
  user study, NASA TLX, cognitive load, task coordination, human-AI interaction,
  software engineering}

%% ── Copyright / DOI stubs (filled in by ACM at acceptance) ─────────────────
%% \acmDOI{10.1145/XXXXXXX.XXXXXXX}
%% \acmISBN{978-x-xxxx-xxxx-x/YY/MM}
%% \acmYear{2026}
%% \copyrightyear{2026}
%% \acmSubmissionID{TOSEM-XXXX-XXXX}

%% ── Additional packages ─────────────────────────────────────────────────────
\usepackage{booktabs}      % professional table rules
\usepackage{multirow}      % multi-row table cells
\usepackage{graphicx}      % figures
\usepackage{microtype}     % improved justification
\usepackage{xspace}        % smart spacing after macros
\usepackage{listings}      % code listings
\usepackage{subcaption}    % subfigures

%% ── Custom macros ───────────────────────────────────────────────────────────
\newcommand{\Weld}{\textsc{Weld}\xspace}
\newcommand{\ContextGate}{\textsc{Context Gate}\xspace}
\newcommand{\CI}{$C_1$\xspace}
\newcommand{\CII}{$C_2$\xspace}
\newcommand{\CIII}{$C_3$\xspace}
\newcommand{\NASATLX}{NASA-TLX\xspace}
\newcommand{\SUS}{SUS\xspace}

%% ── Document body ───────────────────────────────────────────────────────────
\begin{document}

\maketitle

\section{Introduction}
\label{sec:intro}

Modern software development demands constant movement across tools: a task
management board to plan work, a chat platform to coordinate with teammates,
and a code-hosting service to review pull requests and track issues. Each
transition carries a context-switching cost---the cognitive overhead of
re-establishing mental state in a new tool---that fragments attention and slows
progress~\cite{mark2008cost,czerwinski2004diary}. This fragmentation is
compounded when information lives in separate silos: a chat thread that
references a board card requires the developer to manually open two applications
and mentally merge their contents.

Integrated developer workspaces attempt to reduce this overhead by co-locating
surfaces. Yet the empirical record on whether integration alone suffices---and
what further benefit is conferred by embedding an AI agent that can observe and
act across those surfaces---remains thin. Prior work on AI-assisted programming
has focused heavily on code generation~\cite{copilot2022,chen2021codex}, leaving
the broader coordination and project-management layer largely unstudied.

A second open question concerns the role of \emph{context breadth}: how much
workspace state should an agent be allowed to observe? Giving an agent access
to more surfaces may improve its relevance, but it raises concerns about
information overuse, trust, and auditability. Designing studies to isolate this
variable has proven difficult because prior systems enforce context boundaries
in prompts---a mechanism that is opaque, brittle, and unverifiable after the
fact.

\textbf{This paper makes three contributions:}
\begin{enumerate}
  \item \textbf{\Weld system.} We describe \Weld, an open-source all-in-one
    developer workspace (Kanban board, team chat, GitHub integration) augmented
    with a study mode that provisions isolated per-participant workspaces and
    logs fine-grained behavioral telemetry~(\S\ref{sec:system}).

  \item \textbf{Code-enforced context scoping (\ContextGate).} We introduce the
    \ContextGate, a server-side module that assembles agent context from a
    condition-specific allowlist. Because scoping is enforced in code---not
    prompts---the manipulation is mechanically verifiable from audit
    logs~(\S\ref{sec:system}).

  \item \textbf{Empirical study.} We report a between-subjects user study
    ($N$~=~\emph{TBD}) crossing three agent-context conditions against three
    task types, with measures spanning task performance, cognitive load
    (\NASATLX), usability (\SUS), and cross-surface telemetry~(\S\ref{sec:results}).
\end{enumerate}

The remainder of the paper is structured as follows.
Section~\ref{sec:related} surveys related work.
Section~\ref{sec:system} describes \Weld and the study instrumentation.
Section~\ref{sec:study} presents the study design.
Section~\ref{sec:results} reports results.
Section~\ref{sec:discussion} discusses implications, and
Section~\ref{sec:conclusion} concludes.

\section{Related Work}
\label{sec:related}

\subsection{Integrated Developer Workspaces}
\label{sec:related:integrated}

The costs of tool fragmentation in knowledge work have been documented through
diary studies and experience-sampling
methods~\cite{mark2008cost,czerwinski2004diary}. Studies of software engineers
specifically find that inter-tool navigation accounts for a measurable fraction
of working time~\cite{ko2007information,latoza2006maintaining}. Commercial
responses---Linear integrating chat and board, Notion combining docs and tasks,
GitHub's own project boards---have not been systematically evaluated against
single-surface baselines in controlled experiments. We contribute such a
baseline via the \Weld pre-agent phase.

\subsection{AI Agents in Productivity and Coordination Tools}
\label{sec:related:agents}

Large language model (LLM) assistants have demonstrated productivity gains in
code completion~\cite{chen2021codex}, text drafting~\cite{noy2023}, and
information retrieval tasks. Fewer studies have examined agents that both
\emph{read} multi-surface context and \emph{act} on it---for example, moving a
board card \emph{and} drafting a chat post as part of a single agent turn.
Recent work on agentic systems~\cite{yao2023react,shinn2023reflexion} has
largely been evaluated in synthetic benchmarks rather than realistic
collaborative workspaces with human participants.

\subsection{Context-Aware AI Systems}
\label{sec:related:context}

Context-aware systems have long been studied in ubiquitous
computing~\cite{dey2001context} and, more recently, in the
retrieval-augmented generation (RAG) paradigm~\cite{lewis2020rag}. The question
of \emph{how much} context to surface to a language model---the context-breadth
design variable manipulated in this study---has been discussed primarily from an
accuracy perspective. The governance and auditability of context scoping has
received less attention. Our \ContextGate mechanism contributes a design pattern
for making context-scoping decisions both explicit and verifiable, building on
principles of privacy-by-design~\cite{cavoukian2009privacybydesign}.

\section{System: \Weld}
\label{sec:system}

\subsection{Weld Overview}
\label{sec:system:overview}

\Weld is an open-source, full-stack developer workspace built with SvelteKit
(web frontend) and Rust/Axum (backend), using PostgreSQL for persistence and
Redis for real-time pub/sub. It provides three primary collaboration surfaces
that correspond directly to the agent-context conditions in the study:

\begin{itemize}
  \item \textbf{Kanban Board.} Drag-and-drop columns and cards with fractional
    indexing, real-time cross-client synchronization, labels, priorities, and
    assignees.

  \item \textbf{Team Chat.} Channel-based messaging with threading, Markdown
    support, @mentions, and \texttt{[[card:id]]} cross-reference tokens that
    link directly to board cards.

  \item \textbf{GitHub Integration.} GitHub App installation with two-way
    synchronization between board cards and GitHub issues/pull requests,
    webhook-driven updates, and per-participant fixture data for study isolation.
\end{itemize}

\subsection{Study Mode and Workspace Isolation}
\label{sec:system:study}

\Weld includes a \emph{study mode} that provisions each participant with a
fully isolated workspace cloned from a seeded ``PetTracker'' project containing
realistic pre-seeded chat history, board cards, and GitHub fixtures. Isolation
is enforced at the database level: every query in study mode is scoped to the
participant's own \texttt{team\_id} and \texttt{project\_id}, preventing
cross-participant data leakage.

Participants authenticate via single-use access codes rather than account
registration. A resume mechanism allows interrupted sessions to continue from
their saved point across devices.

\subsection{The Context Gate}
\label{sec:system:gate}

The central technical contribution enabling the study is the \ContextGate---a
server-side Rust module that assembles the agent's context before each LLM
call. The gate enforces a \emph{default-deny} policy: it maintains a
per-condition allowlist of context sources, and only queries and includes
sources on the allowlist.

\begin{itemize}
  \item \textbf{\CI} (chat only): the gate queries \emph{only} the team's
    message history. Board and GitHub APIs are never called.
  \item \textbf{\CII} (chat + board): the gate queries messages and board
    cards. GitHub APIs are never called.
  \item \textbf{\CIII} (chat + board + GitHub): the gate queries all three
    sources.
\end{itemize}

Because scoping is a hard code boundary---not a prompt instruction---the
manipulation cannot be subverted by the model and is verifiable from the
\texttt{agent\_context\_sent} audit events logged on every turn. The agent
additionally operates through a double-gated tool executor: tools are offered
to the model only for the allowed condition (first gate), and the executor
re-checks permissions before running any tool call (second gate), logging
rejections as \texttt{agent\_tool\_rejected} events.

\begin{figure}[h]
  \centering
  % TODO: Replace with architecture diagram
  \fbox{\rule{0pt}{4cm}\rule{0.8\linewidth}{0pt}}
  \caption{The \ContextGate architecture. Each LLM call passes through the gate,
    which assembles context strictly from the condition's allowlist and emits a
    \texttt{agent\_context\_sent} audit event. The double-gated tool executor
    independently re-checks permissions before execution.}
  \label{fig:gate}
\end{figure}

\subsection{Telemetry Infrastructure}
\label{sec:system:telemetry}

\Weld logs approximately 40 categories of behavioral events
(Table~\ref{tab:events}), batched from the client and ingested server-side
through an allowlist-only endpoint that rejects any event type not on a fixed
enumeration. This prevents capture of keystrokes, mouse coordinates, or other
data not required by the study protocol.

\begin{table}[h]
  \caption{Selected telemetry event categories.}
  \label{tab:events}
  \begin{tabular}{ll}
    \toprule
    Category & Example event types \\
    \midrule
    Task lifecycle  & \texttt{task\_shown}, \texttt{task\_completed}, \texttt{task\_skipped} \\
    Board actions   & \texttt{card\_moved}, \texttt{card\_priority\_set}, \texttt{card\_edited} \\
    Chat actions    & \texttt{message\_sent}, \texttt{channel\_switched} \\
    Cross-surface   & \texttt{card\_referenced\_from\_chat}, \texttt{pr\_linked\_to\_card} \\
    Agent           & \texttt{agent\_invoked}, \texttt{agent\_tool\_call}, \texttt{agent\_tool\_rejected} \\
    \bottomrule
  \end{tabular}
\end{table}

\section{Study Design}
\label{sec:study}

\subsection{Research Questions}
\label{sec:study:rqs}

\begin{description}
  \item[RQ1] \emph{Integrated workspace value (baseline).}
    Does using an integrated workspace (without an agent) reduce cognitive load
    and time-on-task relative to reported experience with fragmented tooling?

  \item[RQ2] \emph{Effect of context breadth on agent effectiveness.}
    Does increasing the agent's context breadth (\CI $\to$ \CII $\to$ \CIII)
    improve task performance and reduce cognitive load?

  \item[RQ3] \emph{Task-type moderation.}
    Does the effect of context breadth vary across task types (comprehension,
    single-surface, cross-surface)?

  \item[RQ4] \emph{Scaling of benefit with surface integration.}
    Do participants in higher-context conditions show greater cross-surface
    integration behavior (measured by telemetry) and higher integration
    perception scores?
\end{description}

\subsection{Hypotheses}
\label{sec:study:hypotheses}

Based on the research questions and prior work on AI-assisted workflows, we
state four directional hypotheses:

\begin{description}
  \item[H1~(Context breadth and speed).]
    Broader agent context reduces time-on-task: participants in \CIII\ will
    complete agent tasks faster than those in \CII, who will in turn be faster
    than those in \CI.

  \item[H2~(Task-type moderation).]
    The effect of context breadth is moderated by task type. The advantage of
    \CIII\ over \CI\ will be negligible for comprehension and single-surface
    tasks (where cross-surface data is irrelevant) but substantial for
    cross-surface tasks that require coordinating board, chat, and GitHub.

  \item[H3~(Cognitive load).]
    Broader context reduces perceived cognitive load (\NASATLX): \CIII\ $<$
    \CII\ $<$ \CI, driven primarily by effort and mental-demand subscales.

  \item[H4~(Perception--behavior gap).]
    Trust and integration-perception scores will not fully track objective task
    performance: participants may report high trust or high integration
    perception even in conditions where telemetry shows limited cross-surface
    coordination, and vice versa.
\end{description}

\subsection{Participants}
\label{sec:study:participants}

% TODO: Fill in once recruitment is complete.

We target $N$~=~20--30 participants drawn from two complementary populations:
\emph{university CS students} (undergraduate and graduate, recruited via
course mailing lists) and \emph{professional software developers} ($\geq$1 year
of industry experience, recruited via professional networks and online
communities). Recruiting both groups allows us to examine whether the benefit of
context breadth varies with prior experience. Eligibility further required
self-reported familiarity with at least one task-management tool (e.g., Jira,
Linear, GitHub Projects) and one team-chat tool (e.g., Slack, Teams, Discord).
The study was conducted remotely. Participants were compensated
\emph{[amount TBD]}.

\subsection{Conditions}
\label{sec:study:conditions}

Participants were randomly assigned, with balanced block randomization (block
size 3), to one of three between-subjects conditions:

\begin{description}
  \item[\CI~(Chat only, advisory).] The agent observes only team chat history.
    No access to board state or GitHub. Advisory-only: no tools offered to the
    model, so it cannot modify any surface.

  \item[\CII~(Chat + Board, action-capable).] The agent observes chat and the
    full board card state. Available tools: \texttt{create\_card},
    \texttt{move\_card}, \texttt{update\_card}, and \texttt{post\_message}
    (drafts a chat message for the participant to review before sending).

  \item[\CIII~(Chat + Board + GitHub, full action).] The agent observes all
    three surfaces. All \CII\ tools plus \texttt{link\_card\_to\_github},
    enabling it to connect board cards to GitHub issues or pull requests.
\end{description}

\subsection{Tasks}
\label{sec:study:tasks}

The study uses a hybrid design. Each participant completes 8 tasks in fixed
order: two manual \emph{baseline tasks} (C0, no agent available) establishing a
per-person speed and workload anchor, followed by six \emph{agent-assisted
tasks} in their assigned condition. Baseline and agent tasks are
content-disjoint to prevent learning transfer. The agent phase uses the
\texttt{gpt-4o-mini-2024-07-18} model at temperature 0.2 with versioned
condition-specific system prompts pinned for reproducibility.

Task types span both phases:

\begin{description}
  \item[Comprehension.] Synthesize information from one or more surfaces to
    form a judgment (e.g., ``what does the team need most right now?'').
  \item[Single-surface.] Can be completed by acting on a single surface, even
    if cross-surface context helps.
  \item[Cross-surface.] Explicitly requires coordinated action across two or
    more surfaces.
\end{description}

\begin{table}[h]
  \caption{Task design summary. B = Baseline (no agent, C0), T = Agent task.
    Surfaces: Ch = Chat, Bo = Board, GH = GitHub.}
  \label{tab:tasks}
  \begin{tabular}{lllll}
    \toprule
    ID & Phase & Type & Surfaces & Summary \\
    \midrule
    B1 & Baseline & Single-surface & Ch, Bo & Find blocked card from chat; move to In Progress \\
    B2 & Baseline & Cross-surface  & Ch, Bo & Check card status; post reply in chat \\
    \midrule
    T1 & Agent & Comprehension  & Ch, Bo     & Determine next priority; update board \\
    T2 & Agent & Comprehension  & Ch, Bo, GH & Catch-up summary; post in chat \\
    T3 & Agent & Single-surface & Ch         & Draft and post status update in chat \\
    T4 & Agent & Single-surface & Bo         & Set priorities on three backlog cards \\
    T5 & Agent & Cross-surface  & Ch, Bo     & Find bug card from chat; move to In Progress \\
    T6 & Agent & Cross-surface  & Ch, Bo, GH & Link PR to card; announce in chat \\
    \bottomrule
  \end{tabular}
\end{table}

\subsection{Procedure}
\label{sec:study:procedure}

% TODO: Fill in IRB protocol details, session timing.

Each session proceeded as follows:
\begin{enumerate}
  \item Informed consent and eligibility screening.
  \item Orientation to \Weld (un-timed, guided).
  \item \textbf{Baseline phase} (2 tasks, no agent, condition C0): tasks
    presented sequentially; participants self-paced, rated each task before
    advancing.
  \item Phase-transition interstitial: agent panel revealed; \CII/\CIII\
    participants received a brief explanation of available agent actions.
  \item \textbf{Agent phase} (6 tasks): same structure with the agent
    available under the participant's assigned condition (\CI, \CII, or
    \CIII).
  \item Exit survey (\NASATLX, \SUS, integration-perception, AI trust scale,
    open-ended comments, demographics).
  \item Debrief.
\end{enumerate}

Sessions were conducted remotely. Screen recording was \emph{not} collected;
all behavioral data came from in-product telemetry.

\subsection{Measures}
\label{sec:study:measures}

\paragraph{Task performance.}
Task completion (binary, judged from \texttt{task\_final\_state} telemetry) and
time-on-task (seconds between \texttt{task\_shown} and \texttt{task\_completed}
or \texttt{task\_skipped}).

\paragraph{Cognitive load.}
\NASATLX~\cite{hart1988tlx}, six subscales (mental demand, physical demand,
temporal demand, performance, effort, frustration), administered once per phase
in the exit survey.

\paragraph{Usability.}
\SUS~\cite{brooke1996sus}, 10 items, 5-point Likert, yielding a 0--100 score.

\paragraph{Integration perception.}
A five-item scale measuring whether participants perceived workspace surfaces as
connected (e.g., ``I could easily see how chat discussions connected to board
tasks'').

\paragraph{AI trust and effectiveness.}
An eight-item scale (adapted from~\cite{madsen2000measuring}) measuring perceived
agent accuracy, helpfulness, and appropriate scope.

\paragraph{Behavioral telemetry.}
Cross-surface reference events (\texttt{card\_referenced\_from\_chat},
\texttt{pr\_linked\_to\_card}, etc.) as an objective measure of integration
behavior; agent turn counts and tool-call distributions as agent-use indicators.

\section{Results}
\label{sec:results}

% ── PLACEHOLDER ─────────────────────────────────────────────────────────────
% Complete once data collection is finished. This section will report:
%   RQ1 — Baseline cognitive load and task time.
%   RQ2 — Between-condition differences on TLX, SUS, task time, completion.
%   RQ3 — Condition × task-type interaction.
%   RQ4 — Cross-surface telemetry counts and integration-perception scores.

\subsection{Participant Demographics and Data Quality}
\label{sec:results:participants}

% TODO: N per condition, attrition, task completion rates, data exclusions.

\subsection{RQ1: Integrated Workspace Value (Baseline)}
\label{sec:results:rq1}

% TODO: Descriptive statistics and inferential tests for baseline-phase measures.

\begin{table}[h]
  \caption{Baseline-phase NASA-TLX subscale means (SD) by condition. [PLACEHOLDER]}
  \label{tab:baseline_tlx}
  \begin{tabular}{lccc}
    \toprule
    Subscale         & \CI & \CII & \CIII \\
    \midrule
    Mental Demand    & --- & --- & --- \\
    Temporal Demand  & --- & --- & --- \\
    Effort           & --- & --- & --- \\
    Frustration      & --- & --- & --- \\
    \midrule
    Overall Weighted & --- & --- & --- \\
    \bottomrule
  \end{tabular}
\end{table}

\subsection{RQ2: Effect of Context Breadth}
\label{sec:results:rq2}

% TODO: One-way ANOVA or Kruskal-Wallis on TLX, SUS, time-on-task by condition.
%   Post-hoc tests (Bonferroni/Tukey) for pairwise differences.

\begin{figure}[h]
  \centering
  % TODO: Replace with actual figure
  \fbox{\rule{0pt}{4cm}\rule{0.8\linewidth}{0pt}}
  \caption{Task completion time (seconds) by condition and task type.
    Error bars show 95\% CI. [PLACEHOLDER]}
  \label{fig:time_by_condition}
\end{figure}

\subsection{RQ3: Task-Type Moderation}
\label{sec:results:rq3}

% TODO: Condition × task-type interaction (mixed ANOVA or linear mixed model).

\begin{figure}[h]
  \centering
  \fbox{\rule{0pt}{4cm}\rule{0.8\linewidth}{0pt}}
  \caption{NASA-TLX score by condition and task type (comprehension /
    single-surface / cross-surface). [PLACEHOLDER]}
  \label{fig:tlx_interaction}
\end{figure}

\subsection{RQ4: Integration Behavior and Perception}
\label{sec:results:rq4}

% TODO: Cross-surface event counts by condition; integration-perception and
%   AI trust scale scores.

\begin{table}[h]
  \caption{Cross-surface telemetry event counts and integration-perception
    scale means (SD) by condition. [PLACEHOLDER]}
  \label{tab:integration}
  \begin{tabular}{lccc}
    \toprule
    Measure                         & \CI & \CII & \CIII \\
    \midrule
    Cross-surface events (median)   & --- & --- & --- \\
    Integration perception (mean)   & --- & --- & --- \\
    AI trust scale (mean)           & --- & --- & --- \\
    \bottomrule
  \end{tabular}
\end{table}

\section{Discussion}
\label{sec:discussion}

% TODO: Interpret results in light of the four RQs once data are available.

\subsection{Context Breadth as a Design Variable}
\label{sec:discussion:breadth}

% TODO: Discuss the C1 → C3 gradient and what results imply about the marginal
%   value of each additional surface in the agent's context. Connect to the
%   information-foraging literature.

\subsection{Task Type as a Moderator}
\label{sec:discussion:tasktype}

% TODO: Discuss whether cross-surface tasks benefit disproportionately from
%   broader context (as hypothesized). When is a narrower context sufficient?

\subsection{Code-Enforced Context Scoping as a Design Principle}
\label{sec:discussion:gate}

The \ContextGate pattern---enforcing context boundaries in server code rather
than model prompts---has implications beyond this study. First, it makes the
manipulation mechanically verifiable: any reviewer can inspect the
\texttt{agent\_context\_sent} audit logs and confirm that a \CI\ agent never
received board data. Second, it suggests a governance model for production AI
agents in collaborative tools, where organizations need auditable controls over
what information agents can observe. Third, it decouples the \emph{scope} of an
agent from its \emph{capability}, allowing researchers and practitioners to vary
context breadth independently of model capability.

\subsection{Limitations}
\label{sec:discussion:limitations}

% TODO: Address:
%   - Controlled lab setting vs. naturalistic use
%   - Seeded PetTracker content may not generalize to all domains
%   - Single session; no longitudinal effects measured
%   - Between-subjects design: no within-person learning-curve data across conditions
%   - Statistical power (dependent on final N)
%   - Generalization to other LLM models (study used a single pinned model)

\section{Conclusion}
\label{sec:conclusion}

We presented \Weld, an integrated developer workspace with a built-in study
mode, and a between-subjects user study ($N$~=~\emph{TBD}) that isolates the
effect of agent context breadth on task performance, cognitive load, and
cross-surface coordination behavior. The study's primary methodological
contribution is the \ContextGate: a server-side, code-enforced context-scoping
mechanism that makes the agent's information diet both configurable and
auditable, enabling a clean between-subjects manipulation without relying on
prompt engineering.

% TODO: Fill in empirical conclusions once data collection is complete.

Our results [to be reported] suggest [summary of key empirical findings].
Together, these findings have implications for the design of AI agents in
collaborative productivity tools, particularly regarding (1) the value of
gradually expanding context rather than granting full information access by
default, and (2) the importance of audit-logging context assembly decisions as a
trust and governance mechanism.

Future work should examine longitudinal effects of integrated AI agents on team
coordination patterns, study generalization to professional teams working on
real projects, and explore adaptive context scoping---where the agent's
information diet expands or contracts based on task characteristics detected at
runtime.

\Weld and the study infrastructure are open-source and available at
\url{https://github.com/weldtool/weld}.


%% ── Acknowledgements ────────────────────────────────────────────────────────
\begin{acks}
The authors thank the study participants for their time. [Additional
acknowledgements to be added before submission.]
\end{acks}

%% ── Bibliography ────────────────────────────────────────────────────────────
\bibliographystyle{ACM-Reference-Format}
\bibliography{refs}

\end{document}
